\section{Only the columns you asked for}
\label{sec:ch04-only-the-columns}
\index{projections|(}

Three statements is a good answer to the wrong question if each of them drags
back every column of every row. A GraphQL client says exactly which fields it
wants, which is more than a REST client ever tells you, and \code{HotChocolate.Data}
exists to spend that information on the query rather than throw it away.

\index{pagination!cursor}%
This is also where pagination arrives, and the two belong together because they
are configured together. A field that returns every product is a bug waiting for
the catalog to grow. Mosaic gets a second entry point:

\begin{minted}[fontsize=\footnotesize]{csharp}
[UseConnection(IncludeTotalCount = true)]
[UseFiltering]
[UseSorting]
public static async Task<PageConnection<Product>> GetBrowseProductsAsync(
    PagingArguments pagingArguments,
    QueryContext<Product>? query,
    CatalogService catalog,
    CancellationToken cancellationToken)
    => await catalog.BrowseProductsAsync(pagingArguments, query, cancellationToken);
\end{minted}

\index{schema!evidence}%
A second entry point, not a replacement, and the reason is about this book
rather than about storefronts. \code{products} is the field
chapters~\ref{ch:hotchocolate-quickly} and~\ref{ch:the-life-of-a-request}
measured and the field this chapter has been measuring; paginating it would
quietly retire two chapters of evidence and make every number in them
unreproducible. In a real service you would replace \code{products}, deprecate
it, and let chapter~\ref{ch:schema-design-that-survives-change}'s machinery
handle the transition. Here the old field earns its keep by staying comparable.

\subsection*{Three parameters the resolver did not have to build}

\index{HotChocolate!PagingArguments}%
\code{PagingArguments} carries \code{first}, \code{after}, \code{last} and
\code{before}, coerced and validated, and registering
\code{AddPagingArguments()} is what puts those arguments on the
field~\autocite{chillicream2026pagination}.

\index{HotChocolate!QueryContext}%
\code{QueryContext<Product>} is the interesting one. It is a record with three
nullable properties: a predicate built from the \code{where} argument, a sort
definition built from \code{order}, and a selector built from the selection set.
In HotChocolate 16 this is the documented route to projections, in place of the
older \code{[UseProjection]} middleware~\autocite{chillicream2026projections}.
The two attributes on the field add the two arguments and do nothing else;
neither of them touches the query. What touches the query is one call:

\begin{minted}[fontsize=\footnotesize]{csharp}
return await db.Products
    .AsNoTracking()
    .With(query, DefaultOrder)
    .ToPageAsync(pagingArguments, cancellationToken);
\end{minted}

\code{With} applies the three in the order filter, sort, project, which is the
order that keeps the query cheap: narrow the rows first, arrange what is left,
and only then decide which columns to carry~\autocite{chillicream2026projections}.

\index{pagination!keyset}%
\code{ToPageAsync} comes from \code{GreenDonut.Data.EntityFramework} and does
keyset pagination: the cursor encodes the sort keys, so the next page is a seek
rather than an offset. That is why the default order is not decoration:

\begin{minted}[fontsize=\footnotesize]{csharp}
private static SortDefinition<Product> DefaultOrder(SortDefinition<Product> sort)
    => sort.IfEmpty(order => order.AddAscending(p => p.Title))
        .AddAscending(p => p.Id);
\end{minted}

\code{IfEmpty} supplies an ordering only when the caller did not ask for one. The
tiebreaker on the identifier is appended whatever the caller asked for, because
a cursor into an order that is not total cannot say where it points. Two
products with the same title and no tiebreaker, and the page boundary between
them is undefined.

\subsection*{What comes out}

\index{projections!generated SQL}%
Ask for three products and one field each:

\begin{minted}[fontsize=\footnotesize]{graphql}
{ browseProducts(first: 3) { nodes { title } } }
\end{minted}

\begin{minted}[fontsize=\footnotesize]{sql}
SELECT p.title
FROM products AS p
ORDER BY p.title, p.id
LIMIT @p
\end{minted}

One column, because the client asked for one field. Add a filter and an order
and both land in the same statement:

\begin{minted}[fontsize=\footnotesize]{graphql}
{
  browseProducts(
    first: 2
    where: { category: { eq: LIGHTING } }
    order: [{ title: DESC }]
  ) { nodes { title category averageRating } }
}
\end{minted}

\begin{minted}[fontsize=\footnotesize]{sql}
SELECT p.id, p.title, p.category
FROM products AS p
WHERE p.category = @p
ORDER BY p.title DESC, p.id
LIMIT @p1
\end{minted}

\index{projections!and DataLoaders}%
Three columns, and the client asked for two of them. \code{category} is there
because the selection set grew by a field that is a column. \code{id} is there
because \code{averageRating}'s resolver declares that it needs the parent's
identifier, and the projection honoured the declaration. Section~\ref{sec:ch04-two-quiet-failures}
is about what happens when it does not.

\code{averageRating} itself is not in that statement at all, and its absence is
the thing to notice. It belongs to the Reviews domain, it is resolved by a
DataLoader, and it arrives in a second statement covering whatever products the
page turned out to contain:

\begin{minted}[fontsize=\footnotesize]{sql}
SELECT r.product_id AS "ProductId", avg(r.rating::double precision) AS "Average"
FROM reviews AS r
WHERE r.product_id = ANY (@productIds)
GROUP BY r.product_id
\end{minted}

Two statements for a filtered, sorted, projected page with a field from another
domain on it. The averages are computed by PostgreSQL, so the review rows
themselves never cross the wire at all. Projection and batching divide the work
along the same seam the folders do.

\subsection*{A total count that costs nothing extra}

\index{pagination!totalCount}%
Selecting \code{totalCount} alongside the page produces this:

\begin{minted}[fontsize=\footnotesize]{sql}
SELECT (
    SELECT count(*)::int
    FROM products AS p0), p.title
FROM products AS p
ORDER BY p.title, p.id
LIMIT @p
\end{minted}

One statement, not two. The count arrives as a correlated subquery evaluated
once, and I am pointing at it because the received wisdom is that a total count
costs a second round trip, and here it does not. It is still work: a count over
an unfiltered table is a scan, and on a catalog of any size that is the
expensive part of the page rather than the page. But it is not an extra trip,
and the difference matters if you are budgeting round trips rather than CPU.

\subsection*{What projection is not}

\index{projections!limits}%
It is worth being clear about the boundary, because the demonstration above is
flattering. Projection reads properties of the entity out of the selection set,
so anything computed by a resolver is beyond it, whether that resolver belongs
to another domain or is sitting on the entity's own type. The second case is the
one that catches people, because a resolver that does nothing but return a
property looks exactly like a property and is not one. That is the subject of
the next section, and it cost me an afternoon.
\index{projections|)}
